\newpage
\section{Prompts for Reasoning Questions}
\subsection{Response Generation}
\label{subsec:resp_generation_reas}
In response generation for reasoning questions, we replace \texttt{\{Question\}} with the actual question and apply the instruction based on its respective question type. For number-in-general questions, only one of the two bullet points will be used depending on the format of the answer. In particular, if the answer has a specific decimal place, we replace \texttt{\{num\_decimal\}} to the actual number of decimal places. This follows the design of MathVista \cite{lu2024mathvista}.

\begin{example}[Insturctions for Text-in-Chart Questions]
\begin{lstlisting}
    {Question}
    * Your final answer must be grounded to some text that is explicitly written and relevant to the question in the chart.
    * If you need to answer multiple terms, separate them with commas.
    * Unless specified in the question (such as answering with a letter), you are required to answer the full names of subplots and/or labels by default.
\end{lstlisting}
\end{example}

\begin{example}[Insturctions for Text-in-General Questions]
\begin{lstlisting}
    {Question}
    * If there are options in the question, your final answer must conform to one of the options.
    * If there are additional instructions in the question, follow them accordingly.
    * If there are neither options nor additional instructions, you are allowed to respond with a short phrase only.
\end{lstlisting}
\end{example}

\begin{example}[Insturctions for Number-in-Chart Questions]
\begin{lstlisting}
    {Question}
    * Your final answer must be grounded to a number that is exlicitly written and relevant to the question in the chart, even if it's an approximate value.
    * You are allowed to extract numbers within some text when needed.
\end{lstlisting}
\end{example}
\begin{example}[Insturctions for Number-in-General Questions]
\begin{lstlisting}
    {Question}
    * Your final answer must be an exact integer.
    (OR)
    * Your final answer must be a number with {num_decimal} decimal places.
\end{lstlisting}
\end{example}

\newpage
\subsection{Grading}
\label{subsec:grading_reas}
In the grading process, we make an API call for each triplet of (question, ground truth, response). For each type of questions, we provide two in-context learning examples before supplying the triplet. In formatting the template, we replace \texttt{<|question|>}, \texttt{<|ground\_truth|>}, \texttt{<|response|>} with their respective values. Note that for the question, we only supply the original question without answer-type-based instructions that are used to generate the model response.
\newpage
\begin{example}[Grading Instructions for Text-in-Chart Questions]
\begin{lstlisting}
You will be given a question, a ground truth answer and a model response. You need to extract the final answer from the model response, compare it with the ground truth answer, and then assign a binary score. Avoid providing explanations in your response. If there is no provided model response, please leave the extracted answer empty and give a score of 0. 

Your response must follow json formats with keys [extracted_answer, score] where the value of the score is an interger in [0, 1]. You must follow the scoring rules:

### Rules ###
* Give a score of 1 if and only if the final answer and the ground truth answer are referring to the same term. It's acceptable to have different grammar or form (e.g., (*$\alpha$*) and alpha; R^2_{t,h,v,m} and R^2_t,h,v,m). It's also acceptable to have different orders of the terms when question asks for multiple terms.
* Give a score of 0 if any term (e.g., ACC+ and ACC; P-101 and P=101) is different between the final answer and the ground truth.

### Example 1 Starts ###
* Question: What is the name of the curve that intersects y=\lambda exactly three times?
* Ground Truth: P56962
* Response: There is only one curve that intersects y=\lambda exactly three times. The name of the curve is written as P55762.

{
    "extracted_answer": "P55762",
    "score": 0
}
### Example 1 Ends ###


### Example 2 Starts ###
* Question: What is the letter of the subplot where all bars are above 35?
* Ground Truth: (b)
* Response: The letter of the subplot where all bars are above 35 is b.

{
    "extracted_answer": "b",
    "score": 1
}
### Example 2 Ends ###

### Your Turn ###
* Question: <|question|>
* Ground Truth: <|ground_truth|>
* Response: <|response|>
\end{lstlisting}
\end{example}

\newpage

\begin{example}[Grading Instructions for Text-in-General Questions]
\begin{lstlisting}
You will be given a question, a ground truth answer and a model response. You need to extract the final answer from the model response, compare it with the ground truth answer, and then assign a binary score. Avoid providing explanations in your response. If there is no provided model response, please leave the extracted answer empty and give a score of 0. 

Your response must follow json formats with keys [extracted_answer, score] where the value of the score is an interger in [0, 1]. You must follow the scoring rules:

### Rules ###
* If there are predefined options in the question:
    * Give a score of 1 if the final answer matches the ground truth answer exactly.
    * Give a score of 0 if the final answer does not match the ground truth answer.
* If there are no predefined options in the question:
    * Give a score of 1 if the final answer shares the same semantic meaning with the ground truth answer (e.g., "increasing then decreasing" and "moving up then down"; "converge" and "move closer together").
    * Give a score of 0 if the final answer shares different semantic meanings from the ground truth answer (e.g., "increasing then decreasing" and "remain constant"; "converge" and "diverge").

### Example 1 Starts ###
* Question: What is the trend of the red curve between t=10 and t=25?
* Ground Truth: increasing then decreasing
* Response: The red curve is increasing between t=10 and t=25.

{
    "extracted_answer": "increasing",
    "score": 0
}
### Example 1 Ends ###

### Example 2 Starts ###
* Question: What is the interval where the blue curve achieves the maximum value among [0, 50], [50, 100], [100, 150], and [150, 200]?
* Ground Truth: [50, 100]
* Response: The interval where the blue curve achieves the maximum value is [50, 100].

{
    "extracted_answer": "[50, 100]",
    "score": 1
}
### Example 2 Ends ###

### Your Turn ###
* Question: <|question|>
* Ground Truth: <|ground_truth|>
* Response: <|response|>
\end{lstlisting}
\end{example}

\begin{example}[Grading Instructions for Number-in-Chart Questions]
\begin{lstlisting}
You will be given a question, a ground truth answer and a model response. You need to extract the final answer from the model response, compare it with the ground truth answer, and then assign a binary score. Avoid providing explanations in your response. If there is no provided model response, please leave the extracted answer empty and give a score of 0. 

Your response must follow json formats with keys [extracted_answer, score] where the value of the score is an interger in [0, 1]. You must follow the scoring rules:

### Rules ###
* Give a score of 1 if and only if the two numbers are exactly equal in values. It's acceptable to have different notations (e.g., 0.01 and 10^-2; 1500 and 1.5e3).
* Give a score of 0 if the two numbers are different in values.

### Example 1 Starts ###
* Question: What is the value of the red curve at t=10?
* Ground Truth: 0.01
* Response: The value of the red curve at t=10 is 0.012.

{
    "extracted_answer": "0.012",
    "score": 0
}
### Example 1 Ends ###

### Example 2 Starts ###
* Question: What is the value of the blue curve at t=50?
* Ground Truth: 1500
* Response: The value of the blue curve at t=50 is 1.5e3.

{
    "extracted_answer": "1.5e3",
    "score": 1
}
### Example 2 Ends ###

### Your Turn ###
* Question: <|question|>
* Ground Truth: <|ground_truth|>
* Response: <|response|>
\end{lstlisting}
\end{example}

\newpage
\begin{example}[Grading Instructions for Number-in-General Questions]
\begin{lstlisting}
You will be given a question, a ground truth answer and a model response. You need to extract the final answer from the model response, compare it with the ground truth answer, and then assign a binary score. Avoid providing explanations in your response. If there is no provided model response, please leave the extracted answer empty and give a score of 0. 

Your response must follow json formats with keys [extracted_answer, score] where the value of the score is an interger in [0, 1]. You must follow the scoring rules:

### Rules ###
* Give a score of 1 if and only if the two numbers are exactly equal in values. It's acceptable to have different notations (e.g., 0.01 and 10^-2; 1500 and 1.5e3).
* Give a score of 0 if the two numbers are different in values.

### Example 1 Starts ###
* Question: What is the value of the red curve at t=10?
* Ground Truth: 0.01
* Response: The value of the red curve at t=10 is 0.012.

{
    "extracted_answer": "0.012",
    "score": 0
}
### Example 1 Ends ###

### Example 2 Starts ###
* Question: What is the value of the blue curve at t=50?
* Ground Truth: 1500
* Response: The value of the blue curve at t=50 is 1.5e3.

{
    "extracted_answer": "1.5e3",
    "score": 1
}
### Example 2 Ends ###

### Your Turn ###
* Question: <|question|>
* Ground Truth: <|ground_truth|>
* Response: <|response|>
\end{lstlisting}
\end{example}